% !TEX root = ../../mixed-sklyanin-cubics.tex
\subsection{全中心係数での実比較核と尺度の決定}
\label{cmo:subsec:all-central-kernel}

三項までの比較から、まず尺度を未定にした実際の同値族を構成する。
その後に、特異係数を含む残余線束圏の smoothness を使って尺度を決める。
判別集合の一致を同値核の存在の代用にはしない。

\begin{theorem}[指定中心族の全 smooth 係数への比較核]
\label{cmo:thm:all-central-commutative-kernel}
$\alpha^{(1)}=(1,4,9)$、$\alpha^{(2)}=(16,25,49)$ と重み $1,100$ を固定し、
\[
 G=J(x^{(1)})+100J(x^{(2)}),\qquad
 L=\alpha^{(1)}\cdot x^{(1)}+100\alpha^{(2)}\cdot x^{(2)},\qquad
 \lambda=\frac{2\sqrt{21}}9
\]
と置く。ここで $J(x,y,z)=(x^3+y^3+z^3+xyz)/3$ である。
可換族を
\begin{equation}\label{cmo:eq:all-central-target}
 X_c=V\bigl(\lambda(G-u^2L)+cu^3\bigr)\subset\mathbb P^6
\end{equation}
とする。このとき次が成立する。
\begin{enumerate}
 \item 全ての $c\in\C$ に対し、残余線束圏の $\C$-線型 dg Morita 同値
 \[
 \mathcal K^{\mathrm{line}}_{A_c}\simeq
 \mathcal K^{\mathrm{line}}_{X_c}
 \]
 を構成できる。これらは $\C[c]$ 上の一つの Morita 双加群と、
 その逆双加群・単位・余単位から得られる。
 \item $U=\mathbb A^1\setminus V(P_{64})$ 上では、この双加群は相対積上 perfect であり、
 全ての $c\in U(\C)$ で実際の同値
 \begin{equation}\label{cmo:eq:all-central-coherent-equivalence}
 \Ku_{\mathrm{coh}}(A_c)\xrightarrow{\ \sim\ }\Ku(X_c)
 \end{equation}
 を与える。この相対 perfect 核と逆核の同値は任意の導来底変更に保たれる。
 \item 境界 $u=0$ の固定した比較核 $E$ に関して、上の同値は
 中心超平面への残余制限と intertwine する。
\end{enumerate}
これは通常の $\C$-線型 dg Morita 比較である。
全 rotation・指定周期・中心作用の標識の同時保存や、
有理数体への核降下は主張しない。
\end{theorem}

\begin{proof}
\paragraph{尺度を未定にした核の構成.}
命題\ref{cmo:prop:central-parameter-comparison}により、
普遍係数上の二項・三項データを、ある定数 $\kappa\ne0$ を用いて
可換族
\[
 \widetilde X_c=V(G-u^2L+\kappa cu^3)
\]
のデータに一致させられる。
この操作は実際の両側の comonad に対して行うので、
任意の二項・三項データが全構造へ延長するという仮定を置かない。
命題\ref{cmo:prop:central-comonad-three-arity}はその比較を全 arity へ延ばし、
系\ref{cmo:cor:central-three-arity-kernel}は coKleisli 代数の比較から
実際の $\C[c]$-線型 Morita 双加群を構成する。
境界制限との intertwiner も同じ構成から得る。
この段階の比較は、特異係数を除かず、proper な残余線束圏上で成立する。

\paragraph{特異係数を含む smoothness の比較.}
非可換側の残余線束圏が smooth でない係数は正確に
$\Sigma=V(P_{64})$ である。
良い係数では定理\ref{cmo:thm:exact-sixty-four}、
悪い係数では定理\ref{cmo:thm:morse-singular-sixty-four}を使う。

可換三次 $X$ では、全 twist が生成する圏は $\operatorname{Perf}(X)$ である。
実際、$\OO_X(1)$ は ample であり、全 twist は $D_{\mathrm{qc}}(X)$ を
compact に生成する。座標の Koszul 列により七つの連続した twist で十分である。
その残余部分が smooth なら、四つの例外線束との有限 gluing も smooth になる。
この場合は coherent 飽和の証明
（主稿の定理\paperref{main}{crs:thm:coherent}）を適用でき、
$\operatorname{Perf}(X)=D^b(\operatorname{coh}X)$ となる。
特異閉点の剰余体は perfect でないので、これは $X$ が滑らかな場合に限る。
逆に $X$ が滑らかなら、その残余圏は smooth proper である。
従って可換側の残余線束圏の非 smooth 集合も、方程式の特異係数集合に等しい。

既に構成した全 $c$ の Morita 同値は smoothness を保存する。
定理\ref{cmo:thm:commutative-discriminant-candidate}により、
$V(G-u^2L+du^3)$ の特異係数集合は $\lambda^{-1}\Sigma$ だから、
\[
 \Sigma=\kappa^{-1}\lambda^{-1}\Sigma,\qquad
 (\kappa\lambda)\Sigma=\Sigma
\]
となる。補題\ref{cmo:lem:critical-spectrum-scaling}は
$\kappa\lambda=\pm1$ を与える。
必要なら $J$ の正規化に用いた平方根の符号を反転する。
これは二項データを保って三項データの符号を反転するので、
$\kappa=\lambda^{-1}$ と選べる。
同じことは可換側で $u\mapsto-u$ を施しても実現できる。
従って $\widetilde X_c$ は式\eqref{cmo:eq:all-central-target}の $X_c$ に等しい。

\paragraph{smooth locus 上の perfect 性と coherent 飽和.}
$U$ 上の可換三次族は smooth projective であり、
その残余圏は相対 smooth proper である。
Morita 同値によって非可換側の相対残余線束圏にもこの性質が移る。
coKleisli の graph 双加群は底上 perfect であり、
両側の相対対角が perfect であるため、相対積上でも perfect となる。
これは系\ref{cmo:cor:central-three-arity-kernel}の相対版で構成した核である。
既存の全 smooth locus の相対模型
（命題\ref{cmo:prop:full-relative-hh}）とも同じ線束生成対象で同定される。

各複素ファイバーでは、定理\ref{cmo:thm:exact-sixty-four}によって
非可換側の線束圏が全 coherent 圏に一致する。
可換側も滑らかなので $\operatorname{Perf}=D^b(\operatorname{coh})$ であり、
式\eqref{cmo:eq:all-central-coherent-equivalence}を得る。
相対 Morita 双加群と逆双加群の単位・余単位は同型として構成されているから、
導来 tensor による任意の底変更後も同型である。
非正則基底上の全 $D^b(\operatorname{coh})$ まで同一視しているわけではない。
\end{proof}

\begin{corollary}[定常 norm の全六十四根の可換 Ku 標的]
\label{cmo:cor:norm-commutative-kernel}
定常 norm $N(c)$ の全六十四根において、
式\eqref{cmo:eq:all-central-coherent-equivalence}の perfect 比較核が存在する。
特に $c=-44828$ について明示可換標的は
\[
 X_{-44828}
   =V\bigl(\lambda(G-u^2L)-44828u^3\bigr)\subset\mathbb P^6
\]
である。この同値を持つ非可換五様体の点スキームは
$E_1\sqcup E_2\sqcup\{p_c\}$ であり、
全ての滑らかな可換三次五様体上の Poisson 零スキームと被約化後にも一致しない。
\end{corollary}
\begin{proof}
$\gcd(N,P_{64})=1$ と系\ref{cmo:cor:fixed-zero-norm-smooth}から全ての norm 根は $U$ に属する。
前定理を適用し、点スキームと Poisson 零集合については
系\ref{ppz:cor:reduced-point-poisson-counterexample}を用いる。
後者の不一致の証明自体は Ku 比較を前提にしていない。
ここでは同じ例に実際の可換 Ku 標的を与えた。
\end{proof}

\begin{corollary}[同じ二ブロックの任意の定数補正]
\label{cmo:cor:all-alpha-central-kernel}
二つの Sklyanin 因子 $S(1,-1,-1)$ と重み $1,100$ を固定する。
任意の $\alpha^{(1)},\alpha^{(2)}\in\C^3$ と $c\in\C$ に対して、
\[
 X_{\alpha,c}
 =V\bigl(\lambda(G-u^2L_\alpha)+cu^3\bigr),\qquad
 L_\alpha=\alpha^{(1)}\cdot x^{(1)}
                  +100\alpha^{(2)}\cdot x^{(2)}
\]
と $A_{\alpha,c}$ の残余線束圏は、実際の dg Morita 同値を持つ。
比較は普遍多項式環 $T$ 上で構成できる。
可換族が滑らかな開集合上では相対 perfect 核となり、
各複素ファイバーで
\[
 \Ku_{\mathrm{coh}}(A_{\alpha,c})\simeq\Ku(X_{\alpha,c})
\]
を与える。この族では残余線束圏の smoothness は
$X_{\alpha,c}$ の滑らかさと同値である。
\end{corollary}
\begin{proof}
二項・三項の比較と $\beta,\gamma,\kappa$ は全 $\operatorname{Spec}T$ 上で構成され、
$\alpha,c$ に依存しない。前定理で一つの指定 $\alpha$ における
全 $c$ の smoothness を使って定めた $\kappa=\lambda^{-1}$ は、
従って普遍族にも同じ値で適用される。
相対版の系\ref{cmo:cor:central-three-arity-kernel}が全 $T$ 上の比較を与える。
滑らかな可換開上での相対 perfect 性は前定理と同じである。
非可換側の各ファイバーは、PBW/AS 正則な周囲代数の正則中心三次商であり、
有限中心性、proper な線束模型と有限生成性の仮定を全 $\alpha$ で満たす。
そこで smooth な残余線束圏に有限例外部分を戻し、
主稿の定理\paperref{main}{crs:thm:coherent}を適用すると、
全 coherent 圏への飽和を得る。
\end{proof}

\begin{remark}[比較の範囲と残る課題]
\label{cmo:rem:all-central-kernel-scope}
本節の核は $c$ の形式近傍だけでなく、普遍係数上の実際の随伴から作っている。
従って、指定族の一般点同定や norm 根への DVR 特殊化を別に仮定する必要はない。
特異係数で比較するのは proper な残余線束圏であり、
全 coherent 残余圏や perfect な対角をそこへ無条件に持ち込まない。

一方、交換子パラメータ $h$ の選択済み一般点核と
本節の固定 $h=1$ の核を一つの可逆な $h$-族へ接続する問題は残る。
他の Sklyanin パラメータ、線形補正、複数の中心変数、
全 intrinsic 非可換三次への拡張も本定理に含まれない。
全行列係数を列挙した一つの最小核や、指定された全標識を保存する核は求めていない。
主張した核の正本は、制限随伴、正規化した全 comonad 比較、
coKleisli 生成代数とその graph 双加群による構成である。
一般的な球面関手・Morita・障害理論は既知の入力であり、
本族へのこの適用の網羅的な新規性確認は未了である。
\end{remark}
